\documentclass[11pt,a4paper]{article}

\usepackage[margin=1in]{geometry}
\usepackage{booktabs}
\usepackage{hyperref}
\usepackage{graphicx}
\usepackage{csquotes}
\usepackage[capitalise]{cleveref}
\usepackage{xcolor}

\title{Do Classical Price-Action Events Carry Directional Information?\\
A Pre-Registered Confluence Study on EUR/USD and GBP/USD}
\author{Fiyin Akano}
\date{June 2026 \\ \small \texttt{finance-research-experiments} (formerly \texttt{confluence-lab})}

\hypersetup{
  colorlinks=true,
  linkcolor=blue!50!black,
  urlcolor=blue!50!black,
  citecolor=blue!50!black,
}

\begin{document}
\maketitle

\begin{abstract}
Discretionary traders often say that when two or three price-action
patterns line up at the same level, the level matters more than when only
one pattern fires. We pre-registered a study to test this claim
head-on. The study screened seventy-six classical price-action events
across four timeframes on EUR/USD between 2015 and 2021, then attempted
to reproduce any survivor on a frozen 2022--2024 EUR/USD slice and on
GBP/USD over the same screen window. An early version of the analysis
returned forty-one survivors out of two hundred and eighty-four cells,
including pairs of events that contradicted each other (every channel
top and every channel bottom alive at the same time). That result is
the visible footprint of a session-volatility confound: ATR-normalised
forward movement at \emph{random} fifteen-minute times varies by a
factor of 3.7 across the 24-hour cycle. The protocol was amended to
match each control draw to the event's hour of day, and the survivor
count dropped from forty-one to five. One of those five
(\texttt{trendline\_liquidity\_sweep\_low}) reproduced on the frozen
EUR/USD slice; one (\texttt{channel\_top\_touch}) reproduced on
GBP/USD. The effects are real but small (+0.05 to +0.35 units of
average true range): useful as gate inputs to a separately validated
trading agent, far too small to trade on alone. The repository never
trades.
\end{abstract}

\section{Introduction}\label{sec:intro}

This section explains in plain language what \enquote{confluence} means
to a trader, states the research question, and previews the result.

\paragraph{What confluence means.} A retail trader looking at a
fifteen-minute EUR/USD chart will mark a few things on the screen: the
top of an ascending channel, a 50\,\% fibonacci retracement line, the
previous day's high. When all three of those meet near the same price,
the trader calls it a \emph{confluence} and treats the level as
high-confidence. The empirical content of \enquote{confluence} is
therefore a comparative claim: a level with three classical patterns
stacked should hold or reverse harder than a random level at the same
time of day.

\paragraph{The research question.} We asked two questions in nested form.
First, does any individual classical price-action event carry
directional information beyond a matched random baseline? This is
\textbf{H1}. Second, do specific cross-timeframe combinations carry
information their parts do not? This is \textbf{H2}, and is conditional
on H1 producing survivors worth combining.

\paragraph{The headline result.} Five of two hundred and eighty-four
cells survived screening once the controls were corrected for time of
day. One of those five reproduced on a frozen EUR/USD slice, one
reproduced on GBP/USD, and all five kept the predicted sign on both
out-of-sample looks. The effects are too small to trade on standalone
but are statistically real, and the survivors are concentrated in
fifteen-minute trendline and channel geometry.

\paragraph{Outline.} \Cref{sec:background} states what was published
before about price-action edge and multiple-hypothesis discipline.
\Cref{sec:method} defines the event dictionary, the outcome metric, the
controls, and the statistical pipeline. \Cref{sec:confound} is a
self-contained walk-through of the broken first analysis and the
amendment that fixed it. \Cref{sec:results} reports per-stage numbers.
\Cref{sec:limits} lists what we cannot claim. \Cref{sec:conclusion}
closes.

\section{Background and prior work}\label{sec:background}

This section sketches the published literature this study sits in,
because most of what discretionary traders call \enquote{confluence} is
not in any peer-reviewed paper.

The community vocabulary used here (zones, order blocks, fair-value
gaps, market-structure breaks, liquidity sweeps, premium/discount,
fibonacci retracements and extensions) is the Inner Circle Trader (ICT)
glossary, taught publicly online for over a decade with no formal
validation. Academic studies of FX price-action are sparser. Sweeney
(1986) found simple filter rules in spot forex outperformed buy-and-hold
net of costs over the 1970s. LeBaron (1999) extended that to moving-
average rules with mixed results. Brock, Lakonishok and LeBaron (1992)
applied a bootstrap test to technical trading rules on the Dow.

The multiple-hypothesis machinery used here comes from Benjamini and
Hochberg (1995), whose false discovery rate procedure is the standard
correction when many cells are screened at once. Bailey et al. (2014)
catalogued how backtests overfit when this correction is skipped, and
their framing of \enquote{minimum backtest length} motivated our split
discipline. Lopez de Prado (2018) wrote down the broader research
hygiene this lab tries to follow.

To our knowledge, no paper screens the full ICT-style event vocabulary
under pre-registered multiplicity control and reports the survivor list
honestly. This study is that screen.

\section{Methodology}\label{sec:method}

This section locks the event dictionary, the outcome metric, the
controls, and the statistical pipeline. The full protocol was
committed to the repository before any data was screened.

\subsection{Event dictionary (Stage 0)}

Eighteen detector modules emit seventy-six event types. The categories
are: market-structure breaks and shifts, supply and demand zones, order
blocks and breaker blocks, fair-value gaps, liquidity pools and sweeps,
trendlines and parallel channels, horizontal support and resistance,
chart patterns (triple tops, head and shoulders, triangles, wedges,
flags, rectangles), fibonacci retracements and extensions, and
candlestick patterns. Each event ships with a pre-registered directional
hypothesis: \enquote{a touch should bounce}, \enquote{a break should
continue}, \enquote{a sweep should reverse}, \enquote{a magnet should
draw price}.

All detectors are causal: an event timestamped at bar $t$ uses only bars
$\leq t$. The dictionary was frozen before any Stage-1 statistic was
computed. Fifty-one unit tests pin the detector contracts.

\subsection{Outcome metric}

For an event at bar $t$ with directional hypothesis $d$, the maximum
favourable excursion (MFE) is the largest price move in direction $d$
over the next $H$ bars, divided by the average true range (ATR) of the
last 14 bars at $t$. Horizons by timeframe: $H = 30$ on D1, $H = 20$ on
H4 and H1, $H = 16$ on M15. A binary \emph{hit} variable is also recorded:
1 if the event reaches +1 ATR before $-1$ ATR, 0 otherwise.

The ATR normalisation is what lets us compare across timeframes. It is
also what made the first version of the analysis explode: see
\Cref{sec:confound}.

\subsection{Controls}

For each event a number of \emph{control} draws is created from the same
data. Each control is a randomly chosen bar with a direction sampled
from the event cell's own direction mix. The outcome metric is computed
identically. The control is the null model: a level with no event
attached.

The original protocol drew controls uniformly across all bars. We
amended this to draw controls only from bars at the same \emph{hour of
day} as the event. The reason is in \Cref{sec:confound}.

\subsection{Statistical pipeline}

A two-thousand-shuffle permutation test computes a one-sided $p$-value
on the difference in mean MFE between events and controls within a
cell. The Benjamini--Hochberg procedure controls the false discovery
rate at 5\,\% across all adequately powered cells in the family,
defined as cells with at least 100 events.

Verdicts are written into a four-tier registry:

\begin{itemize}
  \item \texttt{alive}: positive effect, survived false discovery rate.
  \item \texttt{parked\_weak\_effect}: positive raw $p < 0.05$ but
        failed the corrected threshold. Recorded, not claimed.
  \item \texttt{parked\_insufficient\_n}: fewer than 100 events.
        Statistics recorded for future re-look.
  \item \texttt{dead}: negative or zero effect; closed.
\end{itemize}

\subsection{Data splits}

\begin{table}[h]
\centering
\caption{Data splits used in the study. Each (pair, split) is consumed
once.}
\label{tab:splits}
\begin{tabular}{lll}
\toprule
Split & Range & Use \\
\midrule
Screen & 2015-01-01 -- 2021-12-31 & Stage 1 and Stage 2 selection \\
Confirm & 2022-01-01 -- 2024-12-31 & frozen test of screen survivors only \\
Sealed & 2025-01-01 onwards & untouched \\
Cross-pair & GBP/USD, same protocol & frozen replication arm \\
\bottomrule
\end{tabular}
\end{table}

\section{The session-volatility confound and amendment v2.1}\label{sec:confound}

This section is the worked example. It walks through what the original
analysis returned, why it could not be true, and how the fix was
specified before any new data was looked at.

\subsection{The pathological first run}

The first Stage 1 run produced 41 \texttt{alive} cells out of 284, an
\emph{a priori} implausible survivor rate. The survivor list included
\texttt{channel\_top\_touch} and \texttt{channel\_bottom\_touch} both
significantly positive on M15. A channel top is supposed to fade; a
channel bottom is supposed to bounce. Both cannot be alive at the same
significance level on the same data unless the null is broken. The
same pattern appeared for liquidity sweeps in both directions and for
several pairs of fibonacci levels.

\subsection{The diagnostic}

A short diagnostic script measured the control MFE distribution on M15
hour-by-hour, ignoring event labels. The result was striking. At 19:00
UTC, average control MFE was 1.09 ATR. At 05:00 UTC, it was 4.03 ATR.
That is a factor of 3.7 in the supposed null. The mechanism is
mechanical: ATR(14) on M15 sums fifteen-minute ranges over the last 3.5
hours, so at the boundary between a quiet session and an active one,
the next sixteen bars structurally exceed the ATR window. Events
themselves are not uniformly distributed across the clock (the
\texttt{asia\_high\_sweep} family, for instance, has 45\,\% of its
events at 07:00 UTC), so any event family concentrated in active hours
inherited a fake positive lift, and any family concentrated in quiet
hours inherited a fake negative one.

\subsection{Amendment v2.1}

The protocol was amended once, documented in the repository before any
re-analysis, to match each control draw to the event's hour of day in
addition to its direction. No detector parameters changed. No event
outcomes were re-measured. Only the control-sampling code changed. The
uniform-control run is preserved in \texttt{output/} as the cautionary
record.

\subsection{What the fix did}

\begin{table}[h]
\centering
\caption{Effect of the amendment on the survivor count.}
\label{tab:amendment}
\begin{tabular}{lrr}
\toprule
Run & Alive cells & Notes \\
\midrule
Uniform controls (original) & 41 / 284 & contradictory pairs alive \\
Hour-matched controls (amended) & 5 / 284 & all M15; trendline / channel / fib \\
\bottomrule
\end{tabular}
\end{table}

Five survivors is a believable number for a screen of 284 cells under
5\,\% false discovery rate. The contradictions disappeared.

\section{Results}\label{sec:results}

This section reports the numbers stage by stage, leading with the
headline cell and then the table.

\subsection{Stage 1 on EUR/USD screen (hour-matched controls)}

Five cells alive, 16 parked weak, 77 parked for insufficient sample
size, 186 dead. All five survivors are on M15:

\begin{table}[h]
\centering
\caption{Stage 1 survivors on EUR/USD 2015--2021 under hour-matched
controls. Effect is in units of ATR(14). All five passed the corrected
threshold.}
\label{tab:stage1}
\begin{tabular}{lrrr}
\toprule
Event type & $n$ & Effect (ATR) & $p$ \\
\midrule
\texttt{channel\_top\_touch} & 19{,}793 & +0.075 & 0.0005 \\
\texttt{fib\_50\_tag} & 5{,}134 & +0.152 & 0.0005 \\
\texttt{trendline\_liquidity\_sweep\_low} & 1{,}133 & +0.303 & 0.0005 \\
\texttt{fib\_618\_tag} & 4{,}676 & +0.124 & 0.0010 \\
\texttt{fib\_ext\_1272\_tag} & 2{,}482 & +0.189 & 0.0010 \\
\bottomrule
\end{tabular}
\end{table}

\subsection{Out-of-sample confirmation (EUR/USD 2022--2024)}

\texttt{trendline\_liquidity\_sweep\_low} confirmed at +0.308 ATR
($p=0.0065$) on the frozen confirm window. The effect is essentially
unchanged from the screen window (+0.303 in, +0.308 out), which is the
signature of a stable effect rather than a lucky screen. The other four
survivors shrank by roughly half out of sample, classic winner's-curse
attenuation.

\subsection{Cross-pair replication (GBP/USD 2015--2021)}

\texttt{channel\_top\_touch} replicated on GBP/USD at +0.099 ATR
($p=0.0005$). All five survivors kept a positive sign on this second
out-of-sample look as well. Five out of five sign-consistency on two
independent out-of-sample tests has probability roughly 3\,\% under a
random-sign null.

\subsection{Stage 2 (cross-timeframe combinations)}

The strict pre-registered Stage 2 was empty by construction. All five
survivors are M15 cells, so there is no surviving higher-timeframe cell
to use as context. An exploratory Stage 2 that included parked-weak
H1 cells flagged \texttt{equal\_highs\_pool} as a candidate context for
a future pre-registration but does not constitute a claim.

\section{Discussion}\label{sec:disc}

The two validated cells are both reaction-to-geometry effects on M15.
\texttt{trendline\_liquidity\_sweep\_low} is a wick below an ascending
support line that closes back above it; the rule predicts an extra
roughly +0.3 ATR of upward excursion within four hours.
\texttt{channel\_top\_touch} is the first touch of a projected parallel-
channel upper boundary, and the rule predicts extra downward excursion
of roughly +0.1 ATR on EUR/USD and roughly +0.1 ATR on GBP/USD.

Both numbers are small relative to spread. EUR/USD spread on M15 is
roughly 0.1--0.2 ATR(M15) in active hours, which means a +0.1 ATR signal
is mostly inside the cost band. These signals are inputs, not
strategies. A live agent could use them as gate or exit inputs only after
passing its own independent grid, holdout, and walk-forward chain.

The session-volatility result is itself a general methodological
finding: intraday ATR-normalised metrics on M15 require time-of-day
matched controls. Uniform-time controls quietly impose a false null.

\section{Limitations}\label{sec:limits}

\begin{itemize}
  \item Each detector is one auditable operational definition. A dead
        verdict closes that definition, not the underlying folk concept.
  \item Daily timeframe cells are structurally underpowered at $n \geq
        100$ on seven years; a cross-pair panel design is the route
        forward for D1 claims.
  \item MFE measures opportunity, not net profit and loss; transaction
        costs are not modelled.
  \item Two looks at the screen window. The amendment was specified from
        a diagnostic on the screen window itself, so the post-amendment
        $p$-values are conditional on one analysis revision. The
        confirm and sealed windows were never used in the diagnostic.
  \item Indicator-only confluence and cross-family interactions are
        scheduled as future experiments (E008 and E009 in the
        repository registry) and not run here.
\end{itemize}

\section{Conclusion}\label{sec:conclusion}

The headline is five cells out of two hundred and eighty-four, two of
which independently reproduced under hour-matched controls. The
discretionary claim that confluence \enquote{works} is partially
supported on M15 trendline, channel, and fibonacci geometry. It is not
broadly supported across the full ICT-style dictionary. Standalone M15
strategies on these signals are not viable after spread; informed gate
or exit inputs to a separately validated trading agent are. The
repository never trades and these findings will not be used live until
they pass the agent's own validation chain.

\section*{References}

\begin{enumerate}
  \item Benjamini, Y., and Hochberg, Y. (1995). Controlling the false
        discovery rate. \emph{Journal of the Royal Statistical Society,
        Series B}, 57(1), 289--300.
  \item Bailey, D. H., Borwein, J. M., L\'opez de Prado, M., and Zhu, Q. J.
        (2014). Pseudo-mathematics and financial charlatanism: the
        effects of backtest overfitting on out-of-sample performance.
        \emph{Notices of the American Mathematical Society}, 61(5),
        458--471.
  \item Brock, W., Lakonishok, J., and LeBaron, B. (1992). Simple
        technical trading rules and the stochastic properties of stock
        returns. \emph{Journal of Finance}, 47(5), 1731--1764.
  \item LeBaron, B. (1999). Technical trading rule profitability and
        foreign exchange intervention. \emph{Journal of International
        Economics}, 49(1), 125--143.
  \item L\'opez de Prado, M. (2018). \emph{Advances in Financial Machine
        Learning}. Wiley.
  \item Sweeney, R. J. (1986). Beating the foreign exchange market.
        \emph{Journal of Finance}, 41(1), 163--182.
  \item Inner Circle Trader (ICT). Community-taught price-action
        glossary, public archives accessed 2025--2026.
  \item Project artefacts:
        \texttt{EXPERIMENTS.md}, \texttt{PROTOCOL\_DISCIPLINE.md},
        \texttt{DATA\_LEDGER.md},
        \texttt{experiments/E006\_test\_a\_price\_action/REPORT.md}.
\end{enumerate}

\end{document}
